\chapter{\texorpdfstring{\code{pdlmi}}{pdlmi}: Finite Certificate Assembly}\index{pdlmi@\texttt{pdlmi}}
\label{sec:pdlmi}

\code{pdlmi} stores finite certificates assembled from \code{pdvar} expressions or coefficient-backed \code{pdmat} inequalities. Square Hermitian inequality families use semidefinite constraints; other inequalities use entry-wise scalar constraints; \code{pdvar} equality uses direct coefficient equations. Direct assembly is the default, while cell-local P\'olya elevation, a dimension-aware Putinar selector, FullBox with band-limited Gram support (\code{SparseFullBox}), and the dense full-box preordering (\code{FullBox}) are explicit inequality alternatives. It is the package boundary between Bernstein objects and ordinary YALMIP solver calls.

\begin{center}
\begin{tikzpicture}[node distance=4mm and 4mm]
  \node[dp/card,text width=20mm,minimum height=17mm] (direct)
    {\textbf{Direct}\\coefficient signs};
  \node[dp/card,text width=20mm,minimum height=17mm,right=of direct] (polya)
    {\textbf{P\'olya}\\lossless elevation};
  \node[dp/card,text width=20mm,minimum height=17mm,right=of polya] (putinar)
    {\textbf{Putinar}\\1-D parity;\\multi-D singletons};
  \node[dp/card,text width=22mm,minimum height=17mm,below=of direct,xshift=12mm] (sparse)
    {\textbf{SparseFullBox}\\band-limited Gram support};
  \node[dp/card,text width=20mm,minimum height=17mm,right=of sparse] (box)
    {\textbf{FullBox}\\all generator products};
  \node[dp/panel,fit=(direct)(polya)(putinar)(sparse)(box),inner sep=4mm] (options) {};
  \node[dp/label,above=3pt of options] {choose exactly one};
  \node[dp/known,left=6mm of options,align=center,font=\scriptsize,
    text width=24mm,minimum height=18mm] (cmp)
    {\code{lhs <= rhs}\\stored residual};
  \node[dp/contribution,right=6mm of options,align=center,font=\scriptsize,
    minimum width=15mm,minimum height=15mm] (export) {\code{toYalmip}};
  \node[dp/result,right=3mm of export,align=center,font=\scriptsize,
    text width=18mm,minimum height=15mm] (yal) {YALMIP\\Constraint};
  \draw[dp/arrow] (cmp.east) -- (options.west);
  \draw[dp/arrow] (options.east) -- (export.west);
  \draw[dp/arrow] (export)--(yal);
\end{tikzpicture}

\textit{The five certificate paths are alternatives. Selecting P\'olya, Putinar, SparseFullBox, or FullBox rebuilding replaces, rather than compounds, the previous selection.}
\end{center}

\section{API Map}
\label{sec:pdlmi-lookup}

\begin{tabularx}{\textwidth}{@{}lY@{}}
\toprule
Task & Entry \\
\midrule
Construct wrapper &
\hyperref[sec:pdlmi-construction]{\code{pdlmi}},
\hyperref[sec:pdvar-comparison]{\code{<=}/\code{>=}/\code{==}} \\
Count constraints &
\hyperref[sec:pdlmi-direct]{direct coefficient constraints} \\
Select P\'olya elevation &
\hyperref[sec:pdlmi-polya]{\code{UsePolya}, \code{PolyaDegree}, and \code{applyPolya}} \\
Select Putinar certificate &
\hyperref[sec:pdlmi-putinar]{\code{UsePutinar}, \code{PutinarOrder}, and
\code{applyPutinar}} \\
Select SparseFullBox &
\hyperref[sec:pdlmi-sparse-full-box]{\code{UseSparseFullBoxPreorder}, \code{SparseFullBoxOrder}, \code{BandWidth}, and \code{applySparseFullBoxPreorder}} \\
Select FullBox &
\hyperref[sec:pdlmi-full-box]{\code{UseFullBoxPreorder}, \code{FullBoxOrder},
and \code{applyFullBoxPreorder}} \\
Export to YALMIP &
\hyperref[sec:pdlmi-toyalmip]{\code{toYalmip}} \\
Solve &
\hyperref[sec:pdlmi-solver-use]{solver-facing use},
\hyperref[sec:examples]{Examples} \\
\bottomrule
\end{tabularx}

\section{Inspect Certificate State}
\index{PutinarOrder@\texttt{PutinarOrder}}\index{SparseFullBoxOrder@\texttt{SparseFullBoxOrder}}\index{BandWidth@\texttt{BandWidth}}\index{FullBoxOrder@\texttt{FullBoxOrder}}

All \code{pdlmi} properties have private set access and are readable with dot notation.  \code{Constraints} contains the assembled YALMIP entries; \code{Residual} and \code{Relation} retain the original comparison for rebuilding.  The remaining properties record the selected certificate state.

\begin{tabularx}{\textwidth}{@{}>{\raggedright\arraybackslash}p{0.32\textwidth}Y@{}}
\toprule
Property & Meaning \\
\midrule
\code{Constraints} & Stored finite YALMIP constraint entries. \\
\code{Residual}, \code{Relation} & Original \code{pdvar} residual and one of \code{"<="}, \code{">="}, or \code{"=="}. \\
\code{UsePolya}, \code{PolyaDegree} & P\'olya selector and nonnegative degree increment. \\
\code{UsePutinar}, \code{PutinarOrder} & Putinar selector and absolute Gram order. \\
\code{UseSparseFullBoxPreorder}, \code{SparseFullBoxOrder}, \code{BandWidth} & SparseFullBox selector, absolute Gram order, and band width; intermediate support is implemented by overlapping tensor-window PSD cliques. \\
\code{UseFullBoxPreorder}, \code{FullBoxOrder} & Full-box selector and absolute Gram order. \\
\bottomrule
\end{tabularx}

For example, \code{C.Relation}, \code{C.Residual.Degree}, and \code{C.UseSparseFullBoxPreorder} inspect a wrapper. Apply methods return a rebuilt value; assigning these properties is unsupported. An intermediate SparseFullBox result stores all three sparse properties. A width-one request canonicalizes to actual Direct state, where those fields are \code{false}, zero, and zero; a request with \code{BandWidth >= SparseFullBoxOrder + 1} canonicalizes to actual FullBox state, where only \code{UseFullBoxPreorder}/\code{FullBoxOrder} record the selection.

\needspace{12\baselineskip}
\section{Construct \texorpdfstring{\code{pdlmi}}{pdlmi} Comparisons}
\label{sec:pdlmi-construction}

\syntaxhead
\begin{lstlisting}[style=dpsyntax]
C = pdlmi(expr, relation)
C = pdlmi(expr, "==")
C = pdlmi(expr, relation, "UsePolya")
C = pdlmi(expr, relation, UsePolya=true, PolyaDegree=d)
C = pdlmi(expr, relation, "UsePolya", "PolyaDegree", d)
C = pdlmi(expr, relation, "UsePutinar")
C = pdlmi(expr, relation, UsePutinar=true, PutinarOrder=r)
C = pdlmi(expr, relation, PutinarOrder=r)
C = pdlmi(expr, relation, "UseSparseFullBoxPreorder")
C = pdlmi(expr, relation, UseSparseFullBoxPreorder=true)
C = pdlmi(expr, relation, BandWidth=b)
C = pdlmi(expr, relation, SparseFullBoxOrder=r)
C = pdlmi(expr, relation, UseSparseFullBoxPreorder=true, ...
    BandWidth=b, SparseFullBoxOrder=r)
C = pdlmi(expr, relation, "UseFullBoxPreorder")
C = pdlmi(expr, relation, UseFullBoxPreorder=true)
C = pdlmi(expr, relation, FullBoxOrder=r)
C = pdlmi(expr, relation, UseFullBoxPreorder=true, FullBoxOrder=r)
C = pdlmi(expr, relation, UseFullBoxPreorder=false, FullBoxOrder=r)
C = pdlmi(expr, relation, ValidationMode=mode)
C = lhs <= rhs
C = lhs >= rhs
C = lhs == rhs
\end{lstlisting}
\arghead
\begin{itemize}[leftmargin=*]
  \item \code{expr} is a \code{pdvar} residual, or a coefficient-backed \code{pdmat} residual for \code{"<="} or \code{">="}; \code{relation} is exactly \code{"<="}, \code{">="}, or \code{"=="}. Inequalities are classified globally from the original coefficient family. \code{pdvar} equality is entry-wise direct coefficient assembly and accepts rectangular or nonsymmetric residuals. \code{pdmat} equality is instead the scalar logical comparison \code{A==B} and never creates a \code{pdlmi}.
  \item \code{UsePolya} is a logical scalar option with default \code{false}. Its bare form, \code{"UsePolya"}, and \code{UsePolya=true} without a degree select increment one.
  \item \code{PolyaDegree=d} is a finite nonnegative integer scalar. It elevates the residual by \code{d} in every parameter direction before assembling constraints. Supplying it without \code{UsePolya} enables P\'olya and issues \codeerr{pdlmi}{ImplicitUsePolya}; \code{UsePolya=false} with positive \code{d} raises \codeerr{pdlmi}{ConflictingPolyaOptions}.
  \item \code{UsePutinar} is a logical scalar option with default \code{false}. Its bare form and the explicit value \code{true} select the dimension-appropriate Putinar certificate at the smallest admissible order.
  \item \code{PutinarOrder=r} is a finite nonnegative integer scalar and is an absolute Bernstein Gram order, not a degree increment. For assembled residual degree \(M\), its default and minimum are \(\lfloor M/2\rfloor\) in one parameter and \(\lceil M/2\rceil\) in two or more parameters. Supplying the order without \code{UsePutinar} enables Putinar assembly. Pairing \code{UsePutinar=false} with any explicit order, including zero, raises \codeerr{pdlmi}{ConflictingPutinarOptions}.
  \item \code{UseSparseFullBox}\allowbreak\code{Preorder} is a logical scalar option with default \code{false}. Its bare form and the explicit value \code{true} select SparseFullBox with default \code{BandWidth=2}; \code{BandWidth} or \code{SparseFullBoxOrder} alone also enables this family. Explicit \code{UseSparseFullBox}\allowbreak\code{Preorder=false} conflicts with either sparse parameter.
  \item \code{BandWidth=b} is a finite positive integer band width. Intermediate widths restrict the FullBox Gram support and are realized by overlapping axis-aligned tensor-window PSD clique blocks. \code{SparseFullBoxOrder=r} is a finite nonnegative absolute Gram order whose default and minimum are \(\lfloor M/2\rfloor\) in one parameter and \(\lceil M/2\rceil\) otherwise, where \(M=\code{C.Residual.Degree}\). Width one canonicalizes to actual Direct state after order validation; width at least \(r+1\) canonicalizes to actual FullBox state; intermediate widths select SparseFullBox state.
  \item \code{UseFullBoxPreorder} is a logical scalar option with default \code{false}. Its bare form and the explicit value \code{true} select full-box assembly at the smallest admissible order.
  \item \code{FullBoxOrder=r} is a finite nonnegative integer scalar and is an absolute order, not a degree increment. Supplying it without \code{UseFullBoxPreorder} enables full-box assembly. When FullBox is enabled, the minimum is \(\lfloor M/2\rfloor\) for a one-parameter residual of degree \(M\), and \(\lceil M/2\rceil\) for two or more parameters. Explicitly pairing \code{UseFullBoxPreorder=false} with \code{FullBoxOrder=r} suppresses that automatic selection: \code{C} remains direct and records \code{r} only as inspectable metadata; the order is not applied to assembly.
  \item P\'olya, Putinar, SparseFullBox, and FullBox selection are mutually exclusive and apply only to inequalities. Any recognized certificate option on an equality raises \codeerr{pdlmi}{Unsupported}\allowbreak\code{EqualityCertificate} before width or order validation. Enabling more than one inequality family raises \codeerr{pdlmi}{Conflicting}\allowbreak\code{Relaxations}.
  \item \code{ValidationMode=mode} accepts scalar text \code{"fast"} or \code{"strict"}, case-insensitively, and defaults to \code{"fast"}. It is transient and not stored. Every \code{apply...} call accepts its own trailing \code{ValidationMode}; a constructor selection is not inherited by later calls.
  \item \code{C} is direct coefficient-wise assembly unless one opt-in inequality mode is selected. Comparisons such as \code{lhs >= rhs} are the ordinary user-facing entry point and accept compatible numeric, \code{pdmat}, affine \code{sdpvar}, and \code{pdvar} right-hand sides.
\end{itemize}

The inspectable properties \code{Constraints}, \code{Residual}, and \code{Relation} expose the assembled finite constraints and the original comparison data used for rebuilding. \code{UsePolya}/\code{PolyaDegree}, \code{UsePutinar}/\code{PutinarOrder}, \code{UseSparseFullBoxPreorder}/\code{SparseFullBoxOrder}/\code{BandWidth}, and \code{UseFullBoxPreorder}/\code{FullBoxOrder} report the selected mode and parameters. Their set access is private. Comparison-created direct objects have every selector false and every numeric certificate property zero. The full-box explicit-false combination described above is the exception: it retains \code{FullBoxOrder=r} while \code{UseFullBoxPreorder} remains false and the stored constraints remain direct. Putinar and SparseFullBox deliberately have no corresponding exception.

\examplehead
The direct constructor is equivalent to the comparison path when the relation and options are explicit.

\begin{lstlisting}[style=dpmatlab]
yalmip('clear')
P = pdvar(2, {[0 1]}, "symmetric");
C = pdlmi(P, ">=", UsePolya=false, PolyaDegree=0);
wrapperClass = class(C)
constraintCount = numel(C.Constraints)
usePolya = C.UsePolya
polyaDegree = C.PolyaDegree
usePutinar = C.UsePutinar
putinarOrder = C.PutinarOrder
useFullBox = C.UseFullBoxPreorder
fullBoxOrder = C.FullBoxOrder
useSparseFullBox = C.UseSparseFullBoxPreorder
sparseFullBoxOrder = C.SparseFullBoxOrder
bandWidth = C.BandWidth
\end{lstlisting}

\begin{lstlisting}[style=dpoutput]
wrapperClass = 'pdlmi'
constraintCount = 2
usePolya =
  logical
   0
polyaDegree = 0
usePutinar =
  logical
   0
putinarOrder = 0
useFullBox =
  logical
   0
fullBoxOrder = 0
useSparseFullBox =
  logical
   0
sparseFullBoxOrder = 0
bandWidth = 0
\end{lstlisting}

\begin{center}
\begin{tikzpicture}[>=Latex,node distance=9mm,
  box/.style={dp/neutral,font=\scriptsize,
    align=center,minimum width=33mm,minimum height=11mm},
  result/.style={dp/result,
    font=\scriptsize,align=center,minimum width=37mm,minimum height=11mm}]
  \node[box] (lhs) {\code{lhs >= rhs}\\matrix comparison};
  \node[result,right=of lhs] (res) {\code{Residual = lhs-rhs}\\
    \code{Relation = ">="}};
  \node[result,right=of res] (con) {\code{Constraints}\\one per selected certificate};
  \draw[->,thick,dpblue] (lhs) -- (res);
  \draw[->,thick,dpblue] (res) -- (con);
\end{tikzpicture}
\end{center}

\remarkshead
\begin{enumerate}[leftmargin=*,itemsep=0.55\baselineskip]
  \item A residual outside the supported \code{pdvar}/\code{pdmat} forms raises \codeerr{pdlmi}{InvalidExpression}; a function-only \code{pdmat} raises \codeerr{pdlmi}{MissingCoefficientEvidence}; \code{pdlmi(A,"==")} for \code{pdmat} raises \codeerr{pdlmi}{UnsupportedPdmatEquality}; an unsupported relation raises \codeerr{pdlmi}{InvalidRelation}. \item Malformed, unknown, or duplicate options raise \codeerr{pdlmi}{InvalidOptions}, \codeerr{pdlmi}{UnknownOption}, or \codeerr{pdlmi}{DuplicateOption}; invalid validation modes raise \codeerr{pdlmi}{InvalidValidationMode}; nonlogical selectors use their corresponding \code{InvalidUse...} identifier. \item Malformed P\'olya increments, Putinar orders, sparse orders, widths, and full-box orders raise, respectively, \codeerr{pdlmi}{Invalid}\allowbreak\code{PolyaDegree}, \codeerr{pdlmi}{Invalid}\allowbreak\code{PutinarOrder}, \codeerr{pdlmi}{InvalidSparseFullBoxOrder}, \codeerr{pdlmi}{InvalidBandWidth}, and \codeerr{pdlmi}{Invalid}\allowbreak\code{FullBoxOrder}. An integer below an active Gram-family minimum raises \codeerr{pdlmi}{PutinarOrder}\allowbreak\code{TooLow}, \codeerr{pdlmi}{SparseFullBoxOrderTooLow}, or \codeerr{pdlmi}{FullBoxOrder}\allowbreak\code{TooLow}. \item Explicit false plus a Putinar order raises \codeerr{pdlmi}{ConflictingPutinarOptions}; explicit false plus a sparse order or width raises \codeerr{pdlmi}{ConflictingSparseFullBoxOptions}; explicit false plus a positive P\'olya increment raises \codeerr{pdlmi}{ConflictingPolyaOptions}. The full-box explicit-false form remains direct and retains the supplied order as metadata. \needspace{9\baselineskip} \item An inequality with any nonsquare or non-Hermitian original coefficient is assembled entry-wise and emits \codeerr{pdlmi}{ElementwiseInequality}; it is not rejected. Numeric Hermitian testing accepts equality at the inclusive \(10^{-10}\) infinity-norm tolerance. \item Equality is direct-only. Every \code{apply...} method raises\newline \codeerr{pdlmi}{Unsupported}\allowbreak\code{EqualityCertificate} before validating its supplied increment, width, or order.
\end{enumerate}

\section{Assemble Direct Coefficient Constraints}\index{direct assembly}\index{direct coefficient assembly}
\label{sec:pdlmi-direct}

For a \code{pdvar} residual, including equality behavior, each physical cell, local Bernstein coefficient, and rate-vertex row if present creates one stored YALMIP constraint entry. For an ordinary expression with \(N_c\) physical cells and \(N_b\) Bernstein coefficients per cell, the number of entries is \(N_cN_b\); with \(N_v\) rate vertices it is \(N_cN_bN_v\). A semidefinite-mode entry is one matrix constraint. An entry-wise-mode entry is the column-major vector of scalar matrix entries compared independently. An equality entry is the corresponding direct coefficient equation and may be rectangular. For a coefficient-backed known \code{pdmat} inequality, Direct and P\'olya assembly instead scan the complete numeric coefficient family and store exactly one logical cell containing the global certificate; \code{toYalmip} exports that state as one scalar logical.

The idea is the Bernstein convex-hull property. After orienting the original residual \(\mathcal F\) as the sign-normalized positive target \(S\), suppress the already-stacked rate-row index and write
\[
S^{(\mathbf c)}(\boldsymbol\alpha)=
  \sum_{\mathbf i}B_{\mathbf i}^M(\boldsymbol\alpha)
  C^{(\mathbf c)}[\mathbf i],
\]
where \(M\) is the assembled residual degree. The semidefinite direct certificate requires \(C^{(\mathbf c)}[\mathbf i]\succeq0\) for every physical cell \(\mathbf c\) and local label \(\mathbf i\), repeated independently for every stored rate row. This is sufficient, not necessary.

For an entry-wise inequality, replace each matrix coefficient by \code{C(:)} and impose the oriented scalar sign on every component. The global classification is made before P\'olya elevation or Gram assembly, so every coefficient and every certificate in one wrapper uses the same mode. Putinar, SparseFullBox, and FullBox create one independent scalar Gram certificate for each column-major entry. Direct equality simply imposes \code{C(:) == 0}; it does not use a positivity certificate or an implicit tolerance.

\begin{center}
The diagram specializes to one parameter, degree two, and one ordinary (non-rate) row; therefore \(c\) is a scalar one-based cell index and \(i\in\{0,1,2\}\).

\begin{tikzpicture}[>=Latex,node distance=5mm and 7mm,
  c/.style={dp/neutral,font=\scriptsize,
    minimum width=14mm,minimum height=7mm},
  ok/.style={dp/certificate,
    font=\scriptsize,minimum width=19mm,minimum height=7mm}]
  \node[c] (c0) {\(C^{(c)}[0]\)};
  \node[c,right=of c0] (c1) {\(C^{(c)}[1]\)};
  \node[c,right=of c1] (c2) {\(C^{(c)}[2]\)};
  \node[ok,below=of c0] (p0) {\(\succeq0\)};
  \node[ok,below=of c1] (p1) {\(\succeq0\)};
  \node[ok,below=of c2] (p2) {\(\succeq0\)};
  \draw[->,dpblue] (c0) -- (p0); \draw[->,dpblue] (c1) -- (p1);
  \draw[->,dpblue] (c2) -- (p2);
  \node[below=4pt of p1,font=\scriptsize]
    {repeat independently for each cell and rate row};
\end{tikzpicture}
\end{center}

\syntaxhead
\begin{lstlisting}[style=dpsyntax]
C = pdlmi(expr, relation)
C = lhs <= rhs
C = lhs >= rhs
C = lhs == rhs
\end{lstlisting}

The constructor and comparison forms above create direct coefficient constraints when no inequality certificate is selected. For a \code{pdvar} residual, each stored coefficient is constrained in the oriented positive-semidefinite or entry-wise mode; equality constrains each coefficient entry directly to zero. Its stored-constraint count is therefore the number of physical cells times the number of local Bernstein coefficients, with an additional factor for active rate vertices. As stated above, a coefficient-backed known \code{pdmat} inequality instead stores exactly one logical cell containing the global numeric certificate.

\examplehead
\begin{lstlisting}[style=dpmatlab]
yalmip('clear')
P = pdvar(1, {[0 1]}, Degree=2);
C = P >= 0;

constraintCount = numel(C.Constraints)
usePolya = C.UsePolya
usePutinar = C.UsePutinar
useSparseFullBox = C.UseSparseFullBoxPreorder
useFullBox = C.UseFullBoxPreorder
\end{lstlisting}

\begin{lstlisting}[style=dpoutput]
constraintCount = 3
usePolya =
  logical
   0
usePutinar =
  logical
   0
useSparseFullBox =
  logical
   0
useFullBox =
  logical
   0
\end{lstlisting}

Here the degree-two, one-parameter residual has three local Bernstein coefficients, so direct assembly stores three coefficient constraints. The wrapper remains direct because no certificate selector was enabled; a call such as \code{C.applyPolya(1)} returns a new wrapper with an elevated coefficient family and leaves \code{C} unchanged.

\section{Select P\'olya Assembly with \texorpdfstring{\code{applyPolya}}{applyPolya}}
\index{Polya@P\'olya}\index{applyPolya@\texttt{applyPolya}}\index{P\'olya certificate}\index{P\'olya elevation}
\label{sec:pdlmi-polya}

With \code{UsePolya=true}, \code{pdlmi} first elevates the stored residual by \code{PolyaDegree} in every parameter direction. For a \code{pdvar} residual it then constrains every elevated local coefficient and stored rate row: if the assembled residual degree is \(M\), the count is \(N_c (M+d+1)^\ell\) for ordinary expressions and \(N_c (M+d+1)^\ell N_v\) for rate-row expressions. A coefficient-backed known \code{pdmat} residual instead stores exactly one logical cell containing the global numeric certificate, including after P\'olya elevation. An increment \(d=0\) is valid and has the same coefficient count as Direct. For an entry-wise \code{pdvar} inequality, elevation preserves the global mode and the oriented sign is applied independently to every column-major scalar entry of each elevated coefficient. Equality wrappers reject this method.

\syntaxhead
\begin{lstlisting}[style=dpsyntax]
Cpolya = C.applyPolya()
Cpolya = C.applyPolya(d)
Cpolya = C.applyPolya(..., ValidationMode=mode)
\end{lstlisting}

\code{applyPolya} is a value-class method: it returns a new \code{pdlmi}, leaves \code{C} unchanged, and rebuilds from \code{C.Residual}. Selecting a new increment therefore replaces rather than compounds a previous selection. The no-argument form uses one; an invalid increment raises \code{pdlmi:InvalidPolyaDegree}.

\examplehead
\begin{lstlisting}[style=dpmatlab]
yalmip('clear')
P = pdvar(1, {[0 1]}, Degree=1);
direct = P >= 0;
polya = direct.applyPolya(2);
directDegree = direct.PolyaDegree
usePolya = polya.UsePolya
polyaDegree = polya.PolyaDegree
constraintCount = numel(polya.Constraints)
\end{lstlisting}

\begin{lstlisting}[style=dpoutput]
directDegree = 0
usePolya =
  logical
   1
polyaDegree = 2
constraintCount = 4
\end{lstlisting}

\begin{center}
\begin{tikzpicture}[>=Latex,node distance=8mm,
  box/.style={dp/result,
    font=\scriptsize,align=center,minimum width=37mm,minimum height=11mm}]
  \node[box] (m) {degree \(M\) residual\\fixed polynomial};
  \node[box,right=of m] (e) {elevate by \(d\)\\degree \(M+d\)};
  \node[box,right=of e] (s) {test every elevated\\coefficient sign};
  \draw[->,thick,dpblue] (m) -- (e);
  \draw[->,thick,dpblue] (e) -- (s);
  \node[dp/label,above=5pt of e]
    {\(\prod_s((1-\alpha_s)+\alpha_s)^d=1\)};
\end{tikzpicture}

\textit{P\'olya degree is a uniform representation increment, not a new decision degree, a grid refinement, or a Gram order.}
\end{center}

The mathematical background and the limits of fixed-increment coefficient positivity are developed in \cref{app:certificates}.

\section{Select Putinar Assembly with
\texorpdfstring{\code{applyPutinar}}{applyPutinar}}
\index{Putinar}\index{applyPutinar@\texttt{applyPutinar}}
\index{Putinar certificate}\index{Putinar order}
\label{sec:pdlmi-putinar}

For an entry-wise inequality, the constructions below are applied independently to each scalar matrix entry in MATLAB column-major order.  Gram variables are not shared between entries.  Equality wrappers reject Putinar selection.

The Putinar selector is dimension-aware.  With one parameter it deliberately uses the same exact even/odd Markov--Luk\'acs
\index{Markov--Lukacs@Markov--Luk\'acs} assembly as
\code{applyFullBoxPreorder}.  For selected absolute order \(r\), the even form matches at degree \(2r\),
\[
S=\mathcal Z_r^{\mathsf T}Q_0\mathcal Z_r +\alpha(1-\alpha)
      \mathcal Z_{r-1}^{\mathsf T}Q_1\mathcal Z_{r-1},
\]
where the second block is omitted at \(r=0\).  The odd form matches at degree \(2r+1\),
\[
S=(1-\alpha)\mathcal Z_r^{\mathsf T}Q_L\mathcal Z_r +\alpha\mathcal Z_r^{\mathsf T}Q_U\mathcal Z_r.
\]
All present Gram matrices are positive semidefinite. The one-parameter default and minimum are \(r=\lfloor M/2\rfloor\) for assembled residual degree \(M\).

For two or more parameters, the selector instead uses the fixed-order, box-specific singleton-generator Bernstein--Gram quadratic module. After normalizing the relation so that \(S(\boldsymbol\alpha)\) is positive semidefinite, it represents
\[
S(\boldsymbol\alpha)=S_0(\boldsymbol\alpha) +\sum_{s=1}^{\ell}\alpha_s(1-\alpha_s) S_s(\boldsymbol\alpha).
\]
Only singleton box generators appear in this multivariate branch: products such as \(\alpha_1(1-\alpha_1)\alpha_2(1-\alpha_2)\) belong to the full-box preordering, not this module.

\syntaxhead
\begin{lstlisting}[style=dpsyntax]
Cputinar = C.applyPutinar()
Cputinar = C.applyPutinar(r)
Cputinar = C.applyPutinar(..., ValidationMode=mode)
\end{lstlisting}

Here \code{r} is an absolute Bernstein Gram order.  The no-argument form uses the dimension-dependent minima above; an explicit order must satisfy the same applicable bound.  In the multivariate branch, the unweighted block \(S_0\) uses tensor basis degree \(r\boldsymbol1\), while \(S_s\) uses degree \(r\boldsymbol1-\boldsymbol e_s\).  A nominal block with a negative component is omitted, which occurs for every multiplier at \(r=0\).  The target is degree-elevated without changing its represented polynomial, and the multivariate Gram sum is matched exactly at tensor degree \(2r\).

\examplehead
\begin{lstlisting}[style=dpmatlab]
yalmip('clear')
P = pdvar(1, {[0 1]}, Degree=2);
direct = P >= 0;
putinar = direct.applyPutinar();
directUsesPutinar = direct.UsePutinar
directOrder = direct.PutinarOrder
directConstraintCount = numel(direct.Constraints)
selectedUsesPutinar = putinar.UsePutinar
selectedOrder = putinar.PutinarOrder
selectedConstraintCount = numel(putinar.Constraints)
\end{lstlisting}

\begin{lstlisting}[style=dpoutput]
directUsesPutinar =
  logical
   0
directOrder = 0
directConstraintCount = 3
selectedUsesPutinar =
  logical
   1
selectedOrder = 1
selectedConstraintCount = 5
\end{lstlisting}

For this one-parameter scalar degree-two cell, order one uses the exact even Markov--Luk\'acs form: an unweighted Gram block of dimension two and one generator-weighted block of dimension one, followed by three exact coefficient identities.  The result is identical to \code{direct.applyFullBoxPreorder()}.  Each physical cell and each active rate-vertex row receives independent copies of these PSD blocks and identities.

\remarkshead
For the multivariate backend, the smallest useful worked case is \(\ell=2\).  Let \(g_s=\alpha_s(1-\alpha_s)\).  At absolute order \(r\), Putinar uses exactly the singleton-generator module
\[
  S=S_0+g_1S_1+g_2S_2,
\]
where \(S_0\) has Bernstein basis degree \((r,r)\), \(S_1\) has degree \((r-1,r)\), and \(S_2\) has degree \((r,r-1)\).  Each \(S_j\) is a PSD Gram form.  MATLAB elevates the sign-normalized target and all present terms to tensor degree \((2r,2r)\), then creates one exact coefficient match per local Bernstein label.  These PSD blocks and equalities are assembled independently for every physical cell, active rate row, and scalar matrix entry; they are never shared across those axes.

\begin{center}
The diagram below specializes to the even one-parameter example above.

\begin{tikzpicture}[node distance=7mm and 9mm,
  n/.style={dp/neutral,font=\scriptsize,align=center,minimum width=28mm,
    minimum height=11mm},
  c/.style={dp/certificate,font=\scriptsize,align=center,minimum width=31mm,
    minimum height=11mm},
  r/.style={dp/result,font=\scriptsize,align=center,minimum width=29mm,
    minimum height=11mm}]
  \node[n] (target) {positive target \(S\)\\elevate to degree \(2r\)};
  \node[c,right=of target] (grams) {PSD Gram blocks\\\(S_0\) and singleton \(S_s\)};
  \node[r,right=of grams] (match) {exact Bernstein\\coefficient identities};
  \draw[dp/arrow] (target) -- (grams);
  \draw[dp/arrow] (grams) -- (match);
\end{tikzpicture}

\textit{The method changes the finite certificate, not the stored residual, decision degree, or physical grid.}
\end{center}

\code{applyPutinar} is a value-class method: it returns a new \code{pdlmi}, leaves \code{C} unchanged, and rebuilds from the original \code{Residual} and \code{Relation}. Reapplication replaces any earlier P\'olya, Putinar, SparseFullBox, or FullBox selection. It adds no implicit margin and makes no solver call.

\remarkshead
\begin{enumerate}[leftmargin=*]
  \item Malformed orders raise \codeerr{pdlmi}{InvalidPutinarOrder}; an integer below \(\lfloor M/2\rfloor\) in one parameter, or below \(\lceil M/2\rceil\) in two or more parameters, raises \codeerr{pdlmi}{PutinarOrderTooLow}. \item A malformed selector raises \codeerr{pdlmi}{InvalidUsePutinar}; explicit false plus any supplied order raises \codeerr{pdlmi}{ConflictingPutinarOptions}. \item Simultaneous relaxation families raise \codeerr{pdlmi}{ConflictingRelaxations}; expression, relation, and option validation in \cref{sec:pdlmi-construction} still applies. Global inequality classification is unchanged: the family is semidefinite only when every original coefficient across all cells and rate rows is square Hermitian. Otherwise the complete family uses independent entry-wise scalar Putinar certificates in MATLAB column-major order.
\end{enumerate}

\section[Select SparseFullBox with \texorpdfstring{\code{applySparseFullBoxPreorder}}{applySparseFullBoxPreorder}]{%
Select SparseFullBox with\\[-0.15em] \texorpdfstring{\code{applySparseFullBoxPreorder}}{applySparseFullBoxPreorder}}
\index{SparseFullBox}\index{band-limited Gram support}\index{applySparseFullBoxPreorder@\texttt{applySparseFullBoxPreorder}}\index{BandWidth@\texttt{BandWidth}}
\label{sec:pdlmi-sparse-full-box}

SparseFullBox is FullBox with band-limited Gram support. It uses the same one-parameter parity terms, multivariate box-generator masks, relation orientation, and exact target-coefficient identities as FullBox. The implementation realizes the restricted support by replacing each dense Gram block with free positive-semidefinite clique blocks on overlapping axis-aligned windows of the tensor Bernstein basis. This implementation mechanism changes only allowed Gram entries: the original residual \(\mathcal F\), physical grid, decision degree \(m\), assembled residual degree \(M\), generator weights, and coefficient identities remain unchanged.

\syntaxhead
\begin{lstlisting}[style=dpsyntax]
Csp = C.applySparseFullBoxPreorder()
Csp = C.applySparseFullBoxPreorder(b)
Csp = C.applySparseFullBoxPreorder(b, r)
Csp = C.applySparseFullBoxPreorder(..., ValidationMode=mode)
\end{lstlisting}

\arghead
\code{b} is a finite positive integer band width and defaults to two; for intermediate state it is also the side length of each overlapping implementation window. \code{r} is an optional absolute SparseFullBox Gram order and defaults to \(\lfloor M/2\rfloor\) for a one-parameter residual of degree \(M\), or \(\lceil M/2\rceil\) for two or more parameters. Every explicit order is validated against that same minimum before endpoint normalization.

\descriptionhead
For a FullBox Gram term with tensor basis degree \(\boldsymbol d\), define the complete label box \(\mathcal I_{\boldsymbol d}=\prod_s\{0,\ldots,d_s\}\). SparseFullBox sets \(w_s=\min(b,d_s+1)\) and introduces one independent PSD block for every window
\[
  W_{\boldsymbol t}=\prod_{s=1}^{\ell}\{t_s,\ldots,t_s+w_s-1\},
  \qquad 0\le t_s\le d_s-w_s+1.
\]
Each block contributes only through the basis labels in its own \(W_{\boldsymbol t}\); the weighted Bernstein maps of all windows and all parity/mask terms are summed and matched exactly to every coefficient of the sign-normalized target. In one parameter, \(b=2\) permits adjacent basis-label interactions and gives block-tridiagonal Gram support, while \(b=3\) gives block-pentadiagonal support. Increasing the width gives a nested sufficient hierarchy because a smaller-window PSD block can be embedded in a larger window, but feasibility need not improve strictly on any particular model.

\begin{center}
\begin{tikzpicture}[node distance=7mm and 8mm,
  basis/.style={dp/neutral,font=\scriptsize,align=center,minimum width=13mm,minimum height=8mm},
  win/.style={dp/certificate,font=\scriptsize,align=center,minimum width=31mm,minimum height=10mm},
  result/.style={dp/result,font=\scriptsize,align=center,minimum width=37mm,minimum height=12mm}]
  \node[basis] (z0) {\(0\)};
  \node[basis,right=of z0] (z1) {\(1\)};
  \node[basis,right=of z1] (z2) {\(2\)};
  \node[basis,right=of z2] (z3) {\(3\)};
  \node[win,below=8mm of z0,xshift=10mm] (w0) {window \(\{0,1\}\)\\\(Q_0\succeq0\)};
  \node[win,right=of w0] (w1) {window \(\{1,2\}\)\\\(Q_1\succeq0\)};
  \node[win,right=of w1] (w2) {window \(\{2,3\}\)\\\(Q_2\succeq0\)};
  \node[result,below=of w1] (match) {sum window contributions\\and match every coefficient};
  \draw[dp/arrow] (w0) -- (match);
  \draw[dp/arrow] (w1) -- (match);
  \draw[dp/arrow] (w2) -- (match);
\end{tikzpicture}

\textit{One-dimensional width two uses every adjacent basis window; tensor problems take the Cartesian window in every coordinate rather than adjacency in a flattened label list.}
\end{center}

\examplehead
\begin{lstlisting}[style=dpmatlab]
yalmip('clear')
P = pdvar(1, {[0 1]}, Degree=4);
direct = P >= 0;
sparse = direct.applySparseFullBoxPreorder(2, 2);
directEndpoint = direct.applySparseFullBoxPreorder(1, 2);
denseEndpoint = direct.applySparseFullBoxPreorder(3, 2);
sparseState = [sparse.UseSparseFullBoxPreorder, ...
    sparse.SparseFullBoxOrder, sparse.BandWidth, ...
    sparse.UseFullBoxPreorder, numel(sparse.Constraints)]
directEndpointState = [directEndpoint.UseSparseFullBoxPreorder, ...
    directEndpoint.UseFullBoxPreorder, directEndpoint.BandWidth, ...
    numel(directEndpoint.Constraints)]
denseEndpointState = [denseEndpoint.UseSparseFullBoxPreorder, ...
    denseEndpoint.UseFullBoxPreorder, denseEndpoint.FullBoxOrder, ...
    numel(denseEndpoint.Constraints)]
\end{lstlisting}
\begin{lstlisting}[style=dpoutput]
sparseState = 1×5 double
     1     2     2     0     8
directEndpointState = 1×4 double
     0     0     0     5
denseEndpointState = 1×4 double
     0     1     2     7
\end{lstlisting}

The intermediate result uses three PSD blocks---two width-two windows for the unweighted degree-two basis and one complete width-two window for the generator-weighted degree-one basis---followed by five exact coefficient identities. Width one returns actual Direct state with no sparse metadata or Gram variables. Since \(b=3=r+1\), the dense endpoint returns actual FullBox state with two dense PSD blocks followed by the same five identities. These are canonical public states, not merely equivalent constraint lists.

\remarkshead
\begin{enumerate}[leftmargin=*]
  \item \code{applySparseFullBoxPreorder} is a value-class method: it returns a new wrapper, leaves the source unchanged, and rebuilds from \code{Residual} and \code{Relation}, replacing any earlier P\'olya, Putinar, SparseFullBox, or FullBox selection. It adds no positivity margin and makes no solver call.
  \item Every physical cell and active rate row owns independent window Gram blocks and coefficient identities. Semidefinite mode builds matrix Gram blocks; entry-wise mode builds an independent scalar certificate for every matrix entry in MATLAB column-major order.
  \item Width one canonicalizes to actual Direct state. A width satisfying \code{b >= r + 1} canonicalizes to actual FullBox state. For \(1<b<r+1\), the result records \code{UseSparseFullBoxPreorder=true}, \code{SparseFullBoxOrder=r}, and \code{BandWidth=b}.
  \item Invalid widths raise \codeerr{pdlmi}{InvalidBandWidth}. Malformed orders raise \codeerr{pdlmi}{InvalidSparse}\allowbreak\code{FullBoxOrder}; orders below the dimension-dependent minimum raise \codeerr{pdlmi}{SparseFullBox}\allowbreak\code{OrderTooLow}. More than two positional inputs raise \codeerr{pdlmi}{InvalidApplyOptions}.
  \item Constructor selector conflicts use \codeerr{pdlmi}{ConflictingSparse}\allowbreak\code{FullBoxOptions} or \codeerr{pdlmi}{ConflictingRelaxations}. Equality raises \codeerr{pdlmi}{UnsupportedEquality}\allowbreak\code{Certificate} before width or order validation.
\end{enumerate}

\section[Select FullBox with \texorpdfstring{\code{applyFullBoxPreorder}}{applyFullBoxPreorder}]{%
Select FullBox with\\[-0.15em] \texorpdfstring{\code{applyFullBoxPreorder}}{applyFullBoxPreorder}}
\index{FullBox}\index{full-box preordering}\index{applyFullBoxPreorder@\texttt{applyFullBoxPreorder}}
\label{sec:pdlmi-full-box}

The mathematical construction is a full-box preordering in the forward local coordinates \(\boldsymbol\alpha\in[0,1]^\ell\); the software state is \code{FullBox}. It is a box-specific Bernstein--Gram representation, not a general Putinar certificate, general-domain SOS interface, or full-space SOS representation. It adds no implicit positivity margin. For an entry-wise inequality, every column-major scalar entry receives an independent copy of the selected Gram construction. Equality wrappers reject FullBox selection.

\syntaxhead
\begin{lstlisting}[style=dpsyntax]
Cbox = C.applyFullBoxPreorder()
Cbox = C.applyFullBoxPreorder(r)
Cbox = C.applyFullBoxPreorder(..., ValidationMode=mode)
\end{lstlisting}

The no-argument form selects the smallest admissible absolute order. The explicit form selects \code{r}; it does not add \code{r} to the residual degree. Both forms are value-class selectors: they return a new \code{pdlmi}, leave \code{C} unchanged, and rebuild from \code{C.Residual}. Reapplication therefore replaces an earlier FullBox order or a P\'olya, Putinar, or SparseFullBox selection rather than compounding it.

The idea is to represent the oriented positive target as a sum of PSD Gram forms multiplied by functions that are nonnegative on the local box:
\[
S(\boldsymbol\alpha)=
  \sum_{J\subseteq\{1,\ldots,\ell\}}
  \left(\prod_{s\in J}\alpha_s(1-\alpha_s)\right)
  \mathcal Z_J(\boldsymbol\alpha)^{\mathsf T}
Q_J\mathcal Z_J(\boldsymbol\alpha),
  \qquad Q_J\succeq0.
\]
In one parameter the implementation uses the even/odd Markov--Luk\'acs forms, with minimum absolute order \(\lfloor M/2\rfloor\) for assembled residual degree \(M\). For two or more parameters it uses every subset product of the \(\ell\) box generators and minimum order \(\lceil M/2\rceil\). The selected order fixes the dense Bernstein Gram bases; it is not added to \(M\).

\examplehead
\begin{lstlisting}[style=dpmatlab]
yalmip('clear')
P = pdvar(1, {[0 1]}, Degree=2);
direct = P >= 0;
box = direct.applyFullBoxPreorder();
directConstraintCount = numel(direct.Constraints)
useFullBox = box.UseFullBoxPreorder
fullBoxOrder = box.FullBoxOrder
fullBoxConstraintCount = numel(box.Constraints)
\end{lstlisting}

\begin{lstlisting}[style=dpoutput]
directConstraintCount = 3
useFullBox =
  logical
   1
fullBoxOrder = 1
fullBoxConstraintCount = 5
\end{lstlisting}

At this order the scalar degree-two cell has two PSD Gram constraints followed by three exact Bernstein coefficient identities.

\remarkshead
For \(\ell=2\), the FullBox state differs from Putinar by retaining the product generator. With \(g_s=\alpha_s(1-\alpha_s)\), the full-box preordering has the exact backend form
\[
S=S_0+g_1S_1+g_2S_2+g_1g_2S_{12}.
\]
The Bernstein basis degrees are respectively \((r,r)\), \((r-1,r)\), \((r,r-1)\), and \((r-1,r-1)\); each present term is a PSD Gram block. All four terms are elevated and coefficient-matched at tensor degree \((2r,2r)\).  Thus the additional \(S_{12}\) block is precisely the two-generator distinction from the singleton Putinar module, not a change to the residual, grid, or decision degree.  The implementation builds the PSD blocks and exact equalities independently for each cell, rate row, and scalar matrix entry.

\begin{center}
\begin{tikzpicture}[node distance=8mm and 9mm,
  flow/.style={font=\scriptsize,align=center,text width=31mm,
    minimum height=12mm,inner sep=4pt},
  certificate/.style={flow,dp/certificate}]
  \node[flow,dp/neutral] (target) {cell \(c\) target\\Bernstein coefficients};
  \node[certificate,right=22mm of target,yshift=9mm] (q0)
    {unweighted Gram block\\\(Q_0\succeq0\)};
  \node[certificate,right=22mm of target,yshift=-9mm] (q1)
    {box-weighted block\\\(\alpha(1-\alpha)Q_1\succeq0\)};
  \node[flow,dp/result,right=22mm of q0,yshift=-9mm] (match)
    {exact coefficient\\identities};
  \draw[dp/arrow] ([yshift=3mm]target.east) -- ++(8mm,0) |- (q0.west);
  \draw[dp/arrow] ([yshift=-3mm]target.east) -- ++(8mm,0) |- (q1.west);
  \draw[dp/arrow] (q0.east) -| ([yshift=3mm]match.west);
  \draw[dp/arrow] (q1.east) -| ([yshift=-3mm]match.west);
\end{tikzpicture}
\end{center}

Appendix~\ref{app:certificates} derives the one-dimensional parity forms, explains the multivariate subset masks, and distinguishes this implemented box preordering from the multivariate singleton-generator Putinar branch, full-space SOS, and general polynomial parsers. Assembly is independent for every physical cell and every active rate-vertex row; Gram variables are not shared across either boundary.  For \code{>=}, the method represents the residual.  For \code{<=}, it negates the target before forming the positive-semidefinite representation.  Within each cell and row, the stored constraints are ordered as all PSD Gram blocks followed by exact equalities matching every Bernstein coefficient.  Consequently, the number of stored constraints is not the direct coefficient count.

\remarkshead
\begin{enumerate}[leftmargin=*,nosep]
  \item Malformed orders raise \codeerr{pdlmi}{InvalidFullBoxOrder}; an admissible integer below the minimum raises \codeerr{pdlmi}{FullBoxOrderTooLow}. The constructor rejects malformed selectors with \codeerr{pdlmi}{InvalidUse}\allowbreak\code{FullBox}\allowbreak\code{Preorder}. \item Simultaneous P\'olya, Putinar, SparseFullBox, or FullBox selection raises \codeerr{pdlmi}{Conflicting}\allowbreak\code{Relaxations}; expression, relation, and option validation in \cref{sec:pdlmi-construction} still applies. \item Global inequality classification is unchanged: the family is semidefinite only when every original coefficient across all cells and rate rows is square Hermitian. Otherwise the complete family uses independent entry-wise scalar FullBox certificates in MATLAB column-major order.
\end{enumerate}

\section{Export with \texorpdfstring{\code{toYalmip}}{toYalmip}}
\index{toYalmip@\texttt{toYalmip}}\index{YALMIP}
\label{sec:pdlmi-toyalmip}

\syntaxhead
\begin{lstlisting}[style=dpsyntax]
F = toYalmip(C)
F = C.toYalmip()
\end{lstlisting}

\arghead For a \code{pdvar} residual, \code{toYalmip} concatenates the stored constraint cells into a YALMIP constraint array for Direct, P\'olya, Putinar, SparseFullBox, and FullBox objects. For a coefficient-backed known \code{pdmat} residual, Direct and P\'olya instead return the one stored scalar logical certificate. Its numeric check uses the wrapper's global inequality classification and inclusive absolute tolerance \(10^{-10}\): semidefinite mode symmetrizes each coefficient matrix as \((C+C^\mathsf{T})/2\) and tests its eigenvalues, whereas only element-wise mode tests individual scalar entries. A false result emits warning \codeerr{pdlmi}{InconclusiveCertificate}: failure of this sufficient coefficient certificate does not prove a violation or indefiniteness of the continuous function. Known-data Putinar, SparseFullBox, and FullBox objects still export YALMIP constraints and may contain auxiliary Gram variables. The method has no package-specific options and does not add a solver wrapper, objective, margin, or diagnostic layer.

\examplehead
The function-call and dot-call forms return equivalent YALMIP constraint arrays.

\begin{lstlisting}[style=dpmatlab]
yalmip('clear')
P = pdvar(2, {[0 1]}, "symmetric");
C = P >= 0;
wrapperClass = class(C)
storedConstraintCount = numel(C.Constraints)
F1 = toYalmip(C);
F2 = C.toYalmip();
firstIsConstraint = isa(F1, "lmi") || isa(F1, "constraint")
firstConstraintCount = length(F1)
secondIsConstraint = isa(F2, "lmi") || isa(F2, "constraint")
secondConstraintCount = length(F2)
\end{lstlisting}

\begin{lstlisting}[style=dpoutput]
wrapperClass = 'pdlmi'
storedConstraintCount = 2
firstIsConstraint =
  logical
   1
firstConstraintCount = 2
secondIsConstraint =
  logical
   1
secondConstraintCount = 2
\end{lstlisting}

\begin{center}
\begin{tikzpicture}[>=Latex,node distance=8mm,
  box/.style={dp/result,
    font=\scriptsize,align=center,minimum width=34mm,minimum height=11mm}]
  \node[box] (cells) {\code{C.Constraints}\\cell/rate/label order};
  \node[box,right=18mm of cells] (flat) {YALMIP constraint array\\\code{F=toYalmip(C)}};
  \node[box,right=of flat] (solve) {\code{optimize(F,...)}\\solver-owned status};
  \draw[->,thick,dpblue] (cells) -- node[above,font=\scriptsize]{concatenate} (flat);
  \draw[->,thick,dpblue] (flat) -- (solve);
\end{tikzpicture}
\end{center}

After this boundary, the package is no longer assembling Bernstein data; YALMIP owns solver choice, diagnostics, and optimization status.

\examplehead
Known coefficient evidence can be certified without introducing a solver model on the Direct and P\'olya paths.

\begin{lstlisting}[style=dpmatlab]
yalmip('clear')
A = pdmat([0 1], {{1, 2; 3, 4}}, ...
    Degree=1, RateBounds=[-1 2]);
direct = A >= 0;
directCertificate = toYalmip(direct)
polya = direct.applyPolya(1);
polyaCertificate = toYalmip(polya)
\end{lstlisting}

\begin{lstlisting}[style=dpoutput]
directCertificate = logical
   1

polyaCertificate = logical
   1
\end{lstlisting}

\section{Boundaries and Related APIs}
\code{pdvar}, \code{rhodiff}, \code{applyPolya}, \code{applyPutinar}, \code{applySparseFullBox}\allowbreak\code{Preorder}, \code{applyFullBox}\allowbreak\code{Preorder}, \code{toYalmip}, \code{optimize}.

\subsection{Solve with YALMIP}
\label{sec:pdlmi-solver-use}

\code{pdlmi} does not own solver settings or wrap \code{optimize}. After \code{toYalmip}, use ordinary YALMIP and validate the returned status before retaining an objective or recovered coefficient data:

\begin{lstlisting}[style=dpsyntax]
opts = sdpsettings("solver", solverName, "verbose", 0);
sol = optimize([C1.toYalmip, C2.toYalmip], objective, opts);
assert(sol.problem == 0, sol.info)
objectiveValue = value(objective);
\end{lstlisting}

A result with \code{problem ~= 0} is not an optimized certificate value. In particular, numerical, unknown, aborted, memory, and space-exhaustion statuses are inconclusive about both the original continuous problem and the value of \(\gamma\).
